
\section{Suspended sand settling into a dredged trench: van Rijn's flume}

Delft Hydraulics flume experiments reported by van Rijn \cite{vanRijn1986b}
(Figs.~16--17): a flow $h_0 = 0.39$~m deep at $\bar u_0 = 0.51$~m/s carrying
160~$\mu$m sand in equilibrium (0.04~kg/s/m in total, 0.03 in suspension and
0.01 as bedload; suspended size 130~$\mu$m, $w_s = 0.013$~m/s; $k_s = 0.025$~m)
crosses a trench cut across the 0.5~m wide flume. The flow slows over the
trench, the suspended sand settles onto the downstream side slope and the
floor, and the trench fills and migrates downstream. The bed profiles
measured after 15 hours for side slopes of 1:3, 1:7 and 1:10 (the ``x''
markers of Fig.~16, digitised here) are the data.

\anuga{} runs the approach flow as normal flow on a plane with the Manning
$n$ of $k_s$, held by Dirichlet boundaries, with the trench cut into the
plane 5~m from the inflow, and the inflow carrying the measured suspended
supply. One sand fraction carries the measured settling velocity
(Ferguson--Church equivalent diameter 138~$\mu$m), the de Leeuw et al.
(2020) entrainment law with the Rouse near-bed ratio referenced at $0.1h$,
Wong--Parker bedload with open inflow and outflow, and an evolving bed. As
van Rijn did for his own model (``the constant of Eq.~26 was adjusted
somewhat to give a suspended load of 0.03~kg/sm at the inlet''), the
entrainment constant is calibrated so that the approach flow is in
equilibrium at the measured supply, by short pre-runs on the plane without
the trench; with the de Leeuw law the uncalibrated equilibrium is already
0.0217~kg/s/m under the default closure, so the factor is 1.38. What the trench then tests is the
response of the transport to the non-uniform flow: the settling lag, the
deposition on the downstream slope and floor, and the Exner coupling.

\begin{figure}[h]
\begin{center}
\includegraphics[width=0.95\textwidth]{setup.png}
\end{center}
\caption{The flume of van Rijn's trench tests, here test 3 with 1:10 side slopes. Tests 1 and 2 cut the same depth with 1:3 and 1:7 slopes.}
\end{figure}


The flume is one-dimensional, so the mesh carries two rows of 0.1~m cells
across it, and the 15 hours of bed evolution run under a morphological
acceleration factor of 10: the bed change of every step is multiplied by
ten, the suspension, which adapts over a settling length of about 20~m in
under a minute, is left alone, and the run covers 1.5 hours of flow.
Wong--Parker bedload is scaled to the measured 0.01~kg/s/m at the inlet,
as van Rijn did with his own bedload constant, and the inlet carries that
bedload as a prescribed flux, as the flume was fed at a set rate. The
default zero-gradient import, which equals the inflow cell's own export,
feeds back on itself here: a cell that aggrades under the fixed inflow
stage sees its transport and its import rise, and within a few hours the
reach aggrades without bound. A rigid approach strip cures that but starves
the first erodible cell behind it of bedload, which then scours a 25--30~cm
hole that feeds the trench; the prescribed supply avoids both, and the two
compute modes then agree to 0.1~mm over the 15 hours.

The pass criterion is the bed level after 15 hours at the measured points:
RMS error below 0.03~m (the fill is 0.08--0.10~m deep over 7~m of flume,
and van Rijn's own profile model matched it to about 0.01~m), and the
deepest remaining point within 1.5~m of the measured position. The gentle
1:10 trench, the one a depth-averaged model is expected to get, runs
routinely; the 1:3 and 1:7 trenches, where the flow separates over the lip
and van Rijn's own depth-averaged attempts failed, join it under the long
option. The case runs the default closure, which since the two-layer suspension
[D-5] became the default is a near-bed layer $0.2h$ thick with the log-law
velocity split. The three tests then give 2.45, 2.19 and 1.74~cm RMS. The
instantaneous exchange that was the default before gives 3.9, 3.2 and
2.85: its near-bed concentration responds at once to the slower flow over
the trench, where in the flume the sediment high in the water column takes
time to settle through it, so the trench fills about 25\,\% too fast. The
inflow concentration is set from the measured supply through the
profile-integrated transport $\beta \bar u h c$ rather than $\bar u h c$,
$\beta = 0.71$ here, since the layers move at their own speeds.

\begin{figure}
\begin{center}
\includegraphics[width=0.9\textwidth]{bed_profile_test3.png}
\end{center}
\caption{Test 3, side slopes 1:10: bed level at three times against the
measured profile after 15 hours (times are morphological, the simulated time
times the acceleration factor).}
\end{figure}

\begin{figure}
\begin{center}
\includegraphics[width=0.9\textwidth]{suspended_load_test3.png}
\end{center}
\caption{Test 3: suspended load along the flume at the end of the run,
scaled by the inflow supply.}
\end{figure}

\IfFileExists{bed_profile_test1.png}{
\begin{figure}
\begin{center}
\includegraphics[width=0.9\textwidth]{bed_profile_test1.png}
\end{center}
\caption{Test 1, side slopes 1:3 (long run).}
\end{figure}
}{}

\IfFileExists{bed_profile_test2.png}{
\begin{figure}
\begin{center}
\includegraphics[width=0.9\textwidth]{bed_profile_test2.png}
\end{center}
\caption{Test 2, side slopes 1:7 (long run).}
\end{figure}
}{}


\subsection*{Closure options explored on the 1:10 trench}

The case runs whatever the default closure is, which since [D-5] became
the default is the two-layer suspension with the velocity split, the bold
row below. The options were each run on test 3
(slopes 1:10, 15 hours, morphological factor 10) to see how far the fit can
be taken, all on the same calibrated setup. Each is chosen by a keyword of
\texttt{set\_deposition} (or \texttt{set\_bedload}); nothing else in the
script changes. The score is the RMS error at the measured points after
15~h, the position and depth of the deepest remaining point (measured
6.57~m, 0.079~m), and the bed volume change per unit width over the three
reaches of the figure (m$^2$ per unit width, positive is bed gain): the
downstream side slope (3.5--6~m) and the four metres past the toe. The
trench upstream of the first measured point (0--3.5~m) fills back to the
original bed, +0.41~m$^2$, in the flume and in every run, so it is left out.

\begin{table}[h]
\footnotesize
\setlength{\tabcolsep}{3pt}
\newcolumntype{L}[1]{>{\raggedright\arraybackslash}p{#1}}
\newcolumntype{T}[1]{>{\ttfamily\raggedright\arraybackslash}p{#1}}
\begin{tabular}{L{3.6cm}T{3.6cm}L{1.0cm}L{1.4cm}rr}
\hline
\textrm{option} & \textrm{chosen by} & RMS (cm) & deepest (m) & slope & toe \\
\hline
measured & & & 6.57, 0.079 & +0.17 & $-0.26$ \\
\hline
instantaneous exchange (the old default) & adaptation='none' & 2.85 & 7.10, 0.039 & +0.20 & $-0.15$ \\
Galappatti lag [D-3] & adaptation='armanini' & 3.7 & & & \\
carried near-bed ratio [D-4] & adaptation='carried' & 2.7 & & & \\
two-layer, near-bed layer $a$ [D-5] & adaptation=\newline 'two\_layer' & 2.66 & 5.50, 0.084 & +0.08 & $-0.20$ \\
\quad exchange rate $\times0.5$ & exchange\_factor=0.5 & 2.80 & & & \\
\quad exchange rate $\times2$ & exchange\_factor=2 & 2.56 & & & \\
\quad exchange rate $\times4$ & exchange\_factor=4 & 2.59 & & & \\
\quad near-bed layer $0.05h$ & layer\_fraction=0.05 & 4.24 & & & \\
\quad near-bed layer $0.2h$ & layer\_fraction=0.2 & 2.17 & 5.72, 0.061 & +0.14 & $-0.18$ \\
\textbf{two-layer $0.2h$, velocity split [D-5v]} (the default, and the case) & \textrm{the defaults} & \textbf{1.74} & 6.22, 0.054 & +0.17 & $-0.19$ \\
\quad near-bed layer $0.1h$ & layer\_fraction=0.1 & 3.12 & & & \\
\quad near-bed layer $0.3h$ & layer\_fraction=0.3 & 2.37 & & & \\
\quad exchange rate $\times0.5$ & exchange\_factor=0.5 & 1.66 & 6.30, 0.055 & +0.17 & $-0.20$ \\
\quad bedload power 2.1 (van Rijn 1984) & set\_bedload(m=2.1) & 2.00 & 6.90, 0.050 & +0.19 & $-0.19$ \\
\quad Rouse number $/\,(1+2(w_s/u_*)^2)$ [S-2b] & rouse\_beta='van\_rijn' & 2.07 & 5.60, 0.070 & +0.12 & $-0.20$ \\
\quad Rouse number scaled to 0.5 & rouse\_scale=0.63 & 3.18 & 5.40, 0.083 & +0.07 & $-0.18$ \\
\quad shear overshoot $1+\lambda(h/u)\,du/ds$, $\lambda=5$ & \textrm{experimental, not kept} & 1.77 & 6.40, 0.053 & +0.17 & $-0.19$ \\
\hline
\end{tabular}
\caption{Test 3 (slopes 1:10) after 15 hours under each closure option:
RMS error at the measured points, position and depth of the deepest
point, and the bed volume change (m$^2$ per unit width) on the downstream
slope (3.5--6~m) and past the toe (6--10~m). Indented rows vary one
setting of the row above them. Each keyword is an argument of
\texttt{set\_deposition} unless stated.}
\end{table}

What the sequence shows. The instantaneous exchange deposits as soon as
the flow slows, so the fill sits on the upstream part of the slope and the
deepest point is 0.5~m too far downstream and half as deep. The Galappatti
lag scales the whole exchange down and over-corrects. The two-layer
suspension keeps the sediment high in the column as a state, and with the
near-bed layer at $0.2h$ and each layer advected at its log-law mean
velocity (the near-bed layer at about $0.55\,\bar u$, the upper at
$1.05\,\bar u$) the deposit on the downstream slope matches the measured
volume and the deepest point lands within 0.35~m of the measured one. The
same setting gives 2.45 and 2.19~cm on the 1:3 and 1:7 trenches (3.07 and
3.26 with the two-layer defaults). Every remaining option redistributes
the deposit between the floor and the slope and leaves the scour past the
toe at 0.18--0.20~m$^2$ against the measured 0.26: in the model the trench
stops trapping at about 10~h, when the fill reaches the toe, where the
flume's trench was still 8~cm deep there at 15~h. The bedload is not the
cause: the flow over the trench floor keeps the Shields stress seven times
the threshold, so 70\,\% of the bedload crosses the trench even at the
start and all of it once the upstream part has filled. Van Rijn's
concentration profiles (his Fig.~17) show the flume carrying 45\,\% of its
load above the near-bed layer where the Rouse fit at $Z=0.79$ carries
5\,\%, but flattening the fit lowers the near-bed concentration and with it
the deposit on the slope, because the two-layer exchange sets both from
the one fitted ratio. Resolving the column in more layers, with a mixing
profile, is the change that would separate them; it is what van Rijn's own
model did.

\IfFileExists{trench_velocity_profile.png}{
\begin{figure}
\begin{center}
\includegraphics[width=0.95\textwidth]{trench_velocity_profile.png}
\end{center}
\caption{Test 3: the two-layer suspension with and without the log-law
velocity split, at three near-bed layer thicknesses, against the
instantaneous exchange that preceded it.}
\end{figure}
}{}

\IfFileExists{trench_two_layer_sweep.png}{
\begin{figure}
\begin{center}
\includegraphics[width=0.95\textwidth]{trench_two_layer_sweep.png}
\end{center}
\caption{Test 3: the two-layer knobs, exchange rate and layer thickness,
without the velocity split.}
\end{figure}
}{}

\IfFileExists{trench_rouse_correction.png}{
\begin{figure}
\begin{center}
\includegraphics[width=0.95\textwidth]{trench_rouse_correction.png}
\end{center}
\caption{Test 3: the Rouse-number correction on top of the best two-layer
setting. A flatter profile carries more sediment high in the column but
deposits less on the slope.}
\end{figure}
}{}

\IfFileExists{trench_bedload_diag.png}{
\begin{figure}
\begin{center}
\includegraphics[width=0.95\textwidth]{trench_bedload_diag.png}
\end{center}
\caption{Test 3, best two-layer setting: bedload and suspended transport
along the flume through the run. The bedload is only partly trapped and
the suspended transport past the toe has not recovered its inflow value
within the flume.}
\end{figure}
}{}

\endinput
